\section{Betrag eines Vektors}
\label{sec:vector_algebra_magnitude}

Jede durch einen Vektor beschriebene Verschiebung besitzt eine Länge. Im
Folgenden leiten wir diese Größe zunächst geometrisch her und verallgemeinern
sie anschließend auf Vektoren im $\R^n$.

\vspace{\baselineskip}
\noindent
\begin{minipage}[c]{0.56\linewidth}
  Für den 2D-Vektor in Koordinatendarstellung
  \[
    \vec a=\begin{pmatrix}a_1\\a_2\end{pmatrix}
  \]
  bilden die Komponenten $a_1$ und $a_2$ die Katheten eines rechtwinkligen
  Dreiecks, wobei die Länge des Vektors $|\vec a|$ der Hypotenuse entspricht.
  Der Satz des Pythagoras ergibt damit 
  \begin{ClozeArea}[16]{3}
    \[
      |\vec a|^2=a_1^2+a_2^2
      \quad\Longleftrightarrow\quad
      |\vec a|=\sqrt{a_1^2+a_2^2}.
    \]
  \end{ClozeArea}
\end{minipage}\hfill
\begin{minipage}[c]{0.38\linewidth}
  \centering
  \CachedTikzFigure{content/resources/chapters/vector_algebra/tikz/vector_magnitude.tex}
\end{minipage}\par
\vspace{0.6em}
\noindent
Diese geometrisch motivierte Formel lässt sich auf Vektoren im $\R^n$
verallgemeinern.

\begin{DefinitionBox}{Betrag eines Vektors}{vector_algebra_vector_magnitude}
  Der \emph{Betrag} \EnglishTerm{magnitude} eines Vektors
  $\vec a=\begin{pmatrix}a_1\\\vdots\\a_n\end{pmatrix}\in\R^n$ ist gegeben durch
  \[
    |\vec a| = \Cloze{8cm}{\sqrt{a_1^2+a_2^2+\dots+a_n^2}}
  \]
  und entspricht der Länge des Vektors. \\
  Für den Betrag schreiben wir alternativ auch Buchstabe ohne Vektorpfeil, d.h. $$a:=|\vec a|.$$
\end{DefinitionBox}

\begin{ExampleBox}{Den Betrag von Vektoren bestimmen}{vector_algebra_calculate_magnitude}
  Wir bestimmen den Betrag der folgenden Vektoren:\par\smallskip
  \noindent
  \begin{minipage}[c]{0.54\linewidth}
    \begin{enumerate}[label=\alph*),leftmargin=1.7em,itemsep=1em]
      \item Für $\textcolor{THMBlue}{\vec a}=\begin{pmatrix}3\\4\end{pmatrix}$ gilt
        \[
          |\textcolor{THMBlue}{\vec a}|
          =\Cloze{5cm}{\sqrt{3^2+4^2}=\sqrt{25}=5}.
        \]

      \item Für $\textcolor{olive!80!black}{\vec b}=\begin{pmatrix}1\\1\end{pmatrix}$ gilt
        \[
          |\textcolor{olive!80!black}{\vec b}|
          =\Cloze{5cm}{\sqrt{1^2+1^2}=\sqrt2}.
        \]

      \item Für $\vec c=\begin{pmatrix}1\\1\\1\end{pmatrix}\in\R^3$ gilt
        \[
          |\vec c|
          =\Cloze{5cm}{\sqrt{1^2+1^2+1^2}=\sqrt3}.
        \]

      \item Für $\vec 0=\begin{pmatrix}0\\\vdots\\0\end{pmatrix}\in\R^n$ gilt
        \[
          |\vec 0|
          =\Cloze{5cm}{\sqrt{0^2+\dots+0^2}=0}.
        \]
    \end{enumerate}
    Der Nullvektor ist der einzige Vektor mit Betrag~$0$.
  \end{minipage}\hfill
  \begin{minipage}[c]{0.42\linewidth}
    \centering
    \CachedTikzFigure{content/resources/chapters/vector_algebra/tikz/vector_magnitude_examples_2d.tex}
  \end{minipage}
\end{ExampleBox}

\begin{TheoremBox}{Rechenregeln für den Betrag}{vector_algebra_magnitude_laws}
  Für $\lambda\in\R$ und $\vec a,\vec b\in\R^n$ gilt
  \begin{enumerate}[label=\arabic*.,leftmargin=1.8em,itemsep=0.9em]
    \item \textbf{Nichtnegativität:}
      $|\vec a|\ge0$.

    \item \textbf{Verhalten bei Skalierung:}
      $|\lambda\vec a|=|\lambda|\,|\vec a|$.

    \item \textbf{Dreiecksungleichung}
      \EnglishTerm{triangle inequality}:
      \par\smallskip
      \noindent
      \begin{minipage}[c]{0.62\linewidth}
        \[
          |\vec a+\vec b|\le|\vec a|+|\vec b|.
        \]
        Die Summe der Längen zweier Dreiecksseiten ist mindestens so groß
        wie die Länge der dritten Seite.
      \end{minipage}\hfill
      \begin{minipage}[c]{0.32\linewidth}
        \centering
        \CachedTikzFigure{content/resources/chapters/vector_algebra/tikz/triangle_inequality.tex}
      \end{minipage}
  \end{enumerate}
\end{TheoremBox}

\clearpage
